\documentclass{article}
\usepackage[utf8]{inputenc}
\usepackage[T1]{fontenc}
\usepackage{graphicx}
\usepackage{booktabs}
\usepackage{amsmath}
\usepackage{amssymb}
\usepackage{hyperref}
\usepackage{geometry}
\usepackage{listings}
\usepackage{xcolor}
\usepackage{float}
\usepackage{comment}

\geometry{a4paper, margin=1in}

% Code styling
\lstset{
    basicstyle=\ttfamily\small,
    breaklines=true,
    frame=single,
    backgroundcolor=\color{gray!10},
    keywordstyle=\color{blue},
    commentstyle=\color{green!50!black},
    stringstyle=\color{red}
}

\title{XGBoost Model for Wage De-censoring:\\Technical Documentation}
\author{Felix Kurz}
\date{December 2024}

\begin{document}

\maketitle

\tableofcontents
\newpage

\section{Executive Summary}

This document describes an XGBoost machine learning model developed to predict uncensored wages from censored Beitragsgrundlagen (BG) data. The model is trained on observations where both censored BG data and uncensored Lohnzettel (LZ) data are available, learning to predict true wages from BG characteristics. Once trained, the model can be applied to future BG data where LZ information is unavailable, effectively de-censoring top-coded wages.

\textbf{Sample:} 10,000 randomly sampled workers (penr) from Beitragsgrundlagen data, years 1994--2012.

\section{Data Preparation and Variable Construction}

\subsection{Target Variable}

\subsubsection{dailywage\_final (from LZ data)}

\begin{itemize}
    \item \textbf{Source:} Lohnzettel dataset (\texttt{lz\_dwage})
    \item \textbf{Construction:}
    \begin{enumerate}
        \item Start with daily wage in nominal Euros from LZ
        \item Adjust to 2000 Euros using CPI: $\text{lz\_dwage} \times \text{cpi}_{2000}$
        \item Ensure positive values: $\max(\text{lz\_dwage} \times \text{cpi}_{2000}, 1)$
        \item Log-transform: $\ln(\text{lz\_dwage} \times \text{cpi}_{2000})$
    \end{enumerate}
    \item \textbf{Data merge:} LZ data joined to BG data by (penr, benr, contryear)
    \item \textbf{Note:} This is the ``gold standard'' true wage used for training. For later versions, only BG features from the Social Security data will be available, if current circumstances continue.
\end{itemize}

\subsection{Training Dataset Construction}

\textbf{training\_data}
\begin{itemize}
    \item \textbf{Base:} \texttt{tobit\_data} (all observations with processed BG data, was already created for the tobit implementation)
    \item \textbf{Merge with LZ:} Left join with Lohnzettel on (penr, benr, contryear)
    \item \textbf{Filter:} Keep only observations where \texttt{lz\_dwage\_log} is not missing
    \item \textbf{Result:} Subset where both censored BG and uncensored LZ wages are available
\end{itemize}

\subsection{Feature Variables (Predictors)}

All features are derived exclusively from BG data and administrative records to ensure they are available in production when LZ data is unavailable.

\subsubsection{Individual Characteristics}

\paragraph{age}
\begin{itemize}
    \item \textbf{Source:} Administrative data (Stammdaten)
    \item \textbf{Construction:} $\text{age} = \text{contryear} - \text{birthyear}$
    \item \textbf{Range:} Restricted to 20--60 years during data preparation
    \item \textbf{Purpose:} Captures life-cycle wage patterns
\end{itemize}

\paragraph{birthyear}
\begin{itemize}
    \item \textbf{Source:} Administrative data (Stammdaten)
    \item \textbf{Construction:} Direct from \texttt{stammdaten\_uniqe\_2022}
    \item \textbf{Purpose:} Captures cohort effects beyond age
\end{itemize}

\paragraph{frau (Gender)}
\begin{itemize}
    \item \textbf{Source:} Administrative data (Stammdaten)
    \item \textbf{Construction:} Binary indicator (1 = female, 0 = male) \newline
    Direct from \texttt{stammdaten\_uniqe\_2022}
    \item \textbf{Purpose:} Captures gender wage gaps
\end{itemize}

\subsubsection{Firm Characteristics}

\paragraph{firm\_size}
\begin{itemize}
    \item \textbf{Source:} Derived from BG data
    \item \textbf{Construction:} Count of unique workers per (benr, contryear) combination 
    \item \textbf{Method:} Group by company number in combination with year and count the remainig unique observations
    \item \textbf{Purpose:} Firm size wage premium
\end{itemize}

\paragraph{firm\_size\_square}
\begin{itemize}
    \item \textbf{Source:} Derived from \texttt{firm\_size}
    \item \textbf{Construction:} $\frac{(\text{firm\_size})^2}{100}$
    \item \textbf{Purpose:} Captures non-linear relationship between firm size and wages
    \item \textbf{Note:} Divided by 100 for numerical scaling
\end{itemize}

\paragraph{oneperson\_dummy}
\begin{itemize}
    \item \textbf{Source:} Derived from \texttt{firm\_size}
    \item \textbf{Construction:} $\mathbb{1}(\text{firm\_size} = 1)$
    \item \textbf{Purpose:} Identifies self-employed or single-person firms
\end{itemize}

\paragraph{smallfirm\_dummy}
\begin{itemize}
    \item \textbf{Source:} Derived from \texttt{firm\_size}
    \item \textbf{Construction:} $\mathbb{1}(\text{firm\_size} < 10)$
\end{itemize}

\subsubsection{Individual Wage History}

\paragraph{mean\_wage}
\begin{itemize}
    \item \textbf{Source:} Individual's wage history from BG data
    \item \textbf{Construction:} Leave-one-out mean to avoid mechanical correlation
    \begin{equation}
        \text{mean\_wage}_i = \begin{cases}
        (\text{mean\_wage}_i - \frac{w_{it}}{n_i}) \times \frac{n_i}{n_i-1} & \text{if } n_i > 1\\
        \text{mean\_wage}_i & \text{if } n_i = 1
        \end{cases}
    \end{equation}
    where $w_{it}$ is the wage for individual $i$ in year $t$, $n_i$ is the number of years observed, and $\text{mean\_wage}_i$ is initially $\frac{1}{n_i}\sum_t w_{it}$.
    \item \textbf{Purpose:} Individual fixed effect, captures persistent wage components
    \item \textbf{Important:} Uses BG wages (censored), not LZ wages
    \item \textbf{Note:} Implementation calculates leave-one-out mean by adjusting the full mean rather than direct exclusion
\end{itemize}

\paragraph{mean\_censored}
\begin{itemize}
    \item \textbf{Source:} Individual's censoring history from BG data
    \item \textbf{Construction:} Leave-one-out mean censoring rate
    \begin{equation}
        \text{mean\_censored}_i = \begin{cases}
            (\text{mean\_censored}_i - \frac{\mathbb{1}(\text{censored}_{it})}{n_i}) \times \frac{n_i}{n_i-1} & \text{if } n_i > 1\\
            \text{mean\_censored}_i & \text{if } n_i = 1
        \end{cases}
    \end{equation}
    where $\text{mean\_censored}_i$ is initially $\frac{1}{n_i}\sum_t \mathbb{1}(\text{censored}_{it})$.
    \item \textbf{Range:} 0 to 1 (proportion of years censored)
    \item \textbf{Purpose:} Indicates individual's tendency to earn above top-coding thresholds
    \item \textbf{Note:} Implementation adjusts the full mean rather than direct exclusion
\end{itemize}

\paragraph{num\_years}
\begin{itemize}
    \item \textbf{Source:} Derived from BG data
    \item \textbf{Construction:} Count of years individual appears in data
    \item \textbf{Purpose:} Experience proxy and data quality indicator
\end{itemize}

\paragraph{oneyear\_dummy}
\begin{itemize}
    \item \textbf{Source:} Derived from \texttt{num\_years}
    \item \textbf{Construction:} $\mathbb{1}(\text{num\_years} = 1)$
    \item \textbf{Purpose:} Flags workers with limited history (less reliable wage estimates)
\end{itemize}

\subsubsection{Coworker Characteristics}

\paragraph{mean\_wage\_coworkers}
\begin{itemize}
    \item \textbf{Source:} Coworkers' wages in same firm-year from BG data
    \item \textbf{Construction:} Leave-one-out mean within (benr, contryear)
    \begin{equation}
        \text{mean\_wage\_coworkers}_{i,f,t} = \begin{cases}
            (\text{mean\_wage\_coworkers}_{f,t} - \frac{w_{it}}{n_{ft}}) \times \frac{n_{ft}}{n_{ft}-1} & \text{if } n_{ft} > 1\\
            \text{mean\_wage\_coworkers}_{f,t} & \text{if } n_{ft} = 1
        \end{cases}
    \end{equation}
    where $f$ is the firm, $t$ is the year, $n_{ft}$ is the number of workers in firm $f$ at year $t$, and $\text{mean\_wage\_coworkers}_{f,t}$ is initially $\frac{1}{n_{ft}}\sum_{j \in f,t} w_{jt}$.
    \item \textbf{Purpose:} Firm wage premium, captures match quality
    \item \textbf{Note:} Uses censored BG wages; implementation adjusts full mean rather than direct exclusion
\end{itemize}

\paragraph{mean\_censored\_coworkers}
\begin{itemize}
    \item \textbf{Source:} Coworkers' censoring status from BG data
    \item \textbf{Construction:} Leave-one-out mean censoring rate in firm-year
    \begin{equation}
        \text{mean\_censored\_coworkers}_{i,f,t} = \begin{cases}
            (\text{mean\_censored\_coworkers}_{f,t} - \frac{\mathbb{1}(\text{censored}_{it})}{n_{ft}}) \times \frac{n_{ft}}{n_{ft}-1} & \text{if } n_{ft} > 1\\
            \text{mean\_censored\_coworkers}_{f,t} & \text{if } n_{ft} = 1
        \end{cases}
    \end{equation}
    where $f$ is the firm, $t$ is the year, $n_{ft}$ is the number of workers in firm $f$ at year $t$, and $\text{mean\_censored\_coworkers}_{f,t}$ is initially $\frac{1}{n_{ft}}\sum_{j \in f,t} \mathbb{1}(\text{censored}_{jt})$.
    \item \textbf{Purpose:} Indicates firm's position in wage distribution
    \item \textbf{Note:} Implementation adjusts full mean rather than direct exclusion
\end{itemize}

\subsubsection{Censoring Information}

\paragraph{interval\_lower}
\begin{itemize}
    \item \textbf{Source:} BG data
    \item \textbf{Construction:} Observed log daily wage from BG (censored)
    \item \textbf{Purpose:} Lower bound of true wage for all observations
\end{itemize}

\paragraph{interval\_upper}
\begin{itemize}
    \item \textbf{Source:} Derived from top-coding thresholds and censoring status
    \item \textbf{Construction: }
    \begin{align}
        \text{topcode\_threshold} &= \text{topcoding}[\text{contryear} - 1971]\\
        \text{daily\_topcode} &= \frac{\text{topcode\_threshold} \times 14}{30}\\
        \text{log\_daily\_topcode} &= \ln(\text{daily\_topcode} \times \text{cpi}_{2000})\\
        \text{interval\_upper} &= \begin{cases}
            \text{log\_daily\_topcode} & \text{if censored} = 1\\
            \text{dailywage\_final} & \text{if censored} = 0
        \end{cases}
    \end{align}
    \item \textbf{Purpose:} Upper bound of true wage, only differs from lower bound for censored observations
\end{itemize}

\paragraph{bg\_censored\_flag}
\begin{itemize}
    \item \textbf{Source:} Derived from multiple censoring rules in BG data
    \item \textbf{Construction:} Multi-step process:
    \begin{enumerate}
        \item Initial censoring based on comparison to annual thresholds:
        \begin{itemize}
            \item If daily wage $\times$ 30 days exceeds threshold
            \item If annual wage exceeds threshold
            \item If base wage component exceeds threshold
        \end{itemize}
        \item Modal daily wage imputation for censored observations
        \item Censoring extended to all wages at or above maximum censored wage
    \end{enumerate}
    \item \textbf{Purpose:} Indicator for whether observation is top-coded
\end{itemize}

\paragraph{bg\_censored\_flag}
\begin{itemize}
    \item \textbf{Source:} Derived from multiple censoring rules in BG data
    \item \textbf{Construction:} Multi-step process:
    \begin{enumerate}
        \item Initial censoring based on comparison to annual thresholds:
        \begin{itemize}
            \item If (bg\_euro + sz\_euro) $\times$ 30 / duration $\geq$ threshold $\times$ 7/6
            \item If (bg\_euro + sz\_euro) $\times$ 30 / dur\_nooverlap $\geq$ threshold $\times$ 7/6  
            \item If (bg\_euro + sz\_euro) $\geq$ threshold $\times$ 14
            \item If bg\_euro $\geq$ threshold $\times$ 12
        \end{itemize}
        \item Modal daily wage imputation for censored observations
        \item Censoring extended to all wages at or above maximum censored wage
    \end{enumerate}
    \item \textbf{Purpose: } Indicator for whether observation is top-coded
\end{itemize}










\subsubsection{Additional Variables Introduced in the Richer V4 Specification}

The V4 model extends the earlier specifications by adding variables on labor market history, nationality, education, job type, fertility, and broader administrative context. These variables are all constructed from administrative data and are available without using LZ wages as inputs.

\paragraph{tenure}
\begin{itemize}
    \item \textbf{Source:} Employment history from BG-linked employment spells
    \item \textbf{Construction:} Within each worker--firm match $(\texttt{penr}, \texttt{benr})$, tenure is defined as the yearly sequence number:
    \begin{equation}
        \text{tenure}_{i,f,t} = 1,2,3,\dots
    \end{equation}
    \item \textbf{Implementation detail:} Tenure is capped at 40:
    \begin{equation}
        \text{tenure} = \min(\text{tenure}, 40)
    \end{equation}
    \item \textbf{Purpose:} Captures match-specific wage growth within the same employer
\end{itemize}

\paragraph{exper}
\begin{itemize}
    \item \textbf{Source:} Worker's observed employment history
    \item \textbf{Construction:} Experience is defined as the within-person sequence of observed yearly labor market records minus one:
    \begin{equation}
        \text{exper}_{i,t} = \text{row number within person} - 1
    \end{equation}
    \item \textbf{Implementation detail:} Experience is capped at 40:
    \begin{equation}
        \text{exper} = \min(\text{exper}, 40)
    \end{equation}
    \item \textbf{Purpose:} Captures accumulated labor market experience over the observed career
\end{itemize}

\paragraph{year\_entry}
\begin{itemize}
    \item \textbf{Source:} Derived from first observed employment year
    \item \textbf{Construction:} First observed year in the panel for each worker:
    \begin{equation}
        \text{year\_entry}_i = \min_t(\text{contryear}_{i,t})
    \end{equation}
    \item \textbf{Purpose:} Captures cohort and labor market entry timing effects
\end{itemize}

\paragraph{age\_entry}
\begin{itemize}
    \item \textbf{Source:} Derived from year of first appearance and birth year
    \item \textbf{Construction:}
    \begin{equation}
        \text{age\_entry}_i = \text{year\_entry}_i - \text{birthyear}_i
    \end{equation}
    \item \textbf{Purpose:} Proxy for education duration and timing of labor market entry
\end{itemize}

\paragraph{return}
\begin{itemize}
    \item \textbf{Source:} Derived from gaps in observed yearly employment records
    \item \textbf{Construction:} Indicator for re-entry after a gap in the observed panel:
    \begin{equation}
        \text{return}_{i,t} = \mathbb{1}(\text{contryear}_{i,t} > \text{contryear}_{i,t-1} + 1)
    \end{equation}
    \item \textbf{Purpose:} Captures interrupted employment careers and re-entry dynamics
\end{itemize}

\paragraph{juj}
\begin{itemize}
    \item \textbf{Source:} Derived from transitions between employment and unemployment spells
    \item \textbf{Construction:} Binary indicator based on adjacent unemployment transitions:
    \begin{equation}
        \text{juj}_{i,t} = \mathbb{1}(\text{job\_unemp}_{i,t}=1 \; \text{or} \; \text{lead(unemp\_job)}_{i,t}=1)
    \end{equation}
    \item \textbf{Purpose:} Captures instability around job-to-unemployment and unemployment-to-job transitions
\end{itemize}

\paragraph{stsl\_missing}
\begin{itemize}
    \item \textbf{Source:} Administrative nationality variable \texttt{stsl}
    \item \textbf{Construction:} Indicator for missing or invalid nationality code:
    \begin{equation}
        \text{stsl\_missing} = \mathbb{1}(\texttt{stsl} \text{ is missing or } \texttt{stsl}=999)
    \end{equation}
    \item \textbf{Purpose:} Distinguishes missing nationality information from observed nationality categories
\end{itemize}

\paragraph{austrian\_dummy}
\begin{itemize}
    \item \textbf{Source:} Administrative nationality variable \texttt{stsl}
    \item \textbf{Construction:}
    \begin{equation}
        \text{austrian\_dummy} = \mathbb{1}(\texttt{stsl}=1)
    \end{equation}
    \item \textbf{Purpose:} Identifies Austrian nationals
\end{itemize}

\paragraph{germany\_dummy}
\begin{itemize}
    \item \textbf{Source:} Administrative nationality variable \texttt{stsl}
    \item \textbf{Construction:}
    \begin{equation}
        \text{germany\_dummy} = \mathbb{1}(\texttt{stsl} \in \{2,12\})
    \end{equation}
    \item \textbf{Purpose:} Separates German nationals from both Austrians and other foreign nationals
\end{itemize}

\paragraph{eu\_dummy}
\begin{itemize}
    \item \textbf{Source:} Administrative nationality variable \texttt{stsl}
    \item \textbf{Construction:} Binary indicator for EU nationality excluding Austria and Germany:
    \begin{equation}
        \text{eu\_dummy} = \mathbb{1}(\texttt{stsl} \in \text{predefined EU code list excluding Austria and Germany})
    \end{equation}
    \item \textbf{Purpose:} Captures broader nationality-related wage differences
\end{itemize}

\paragraph{educ\_cat}
\begin{itemize}
    \item \textbf{Source:} Education imputation file
    \item \textbf{Construction:} Categorical education measure with explicit missing category:
    \begin{equation}
        \text{educ\_cat} =
        \begin{cases}
            0 & \text{if education is missing}\\
            \text{educ} & \text{otherwise}
        \end{cases}
    \end{equation}
    \item \textbf{Purpose:} Captures educational attainment while preserving missingness as information
\end{itemize}

\paragraph{educ\_missing}
\begin{itemize}
    \item \textbf{Source:} Education imputation file
    \item \textbf{Construction:}
    \begin{equation}
        \text{educ\_missing} = \mathbb{1}(\text{education missing})
    \end{equation}
    \item \textbf{Purpose:} Separates missing education information from observed educational categories
\end{itemize}

\paragraph{educ\_imputed}
\begin{itemize}
    \item \textbf{Source:} Education imputation file
    \item \textbf{Construction:} Binary indicator showing whether education was imputed:
    \begin{equation}
        \text{educ\_imputed} = \mathbb{1}(\text{education value was imputed})
    \end{equation}
    \item \textbf{Purpose:} Allows the model to distinguish between directly observed and imputed education values
\end{itemize}

\paragraph{whitecollar\_dummy}
\begin{itemize}
    \item \textbf{Source:} Employment subcategory \texttt{subcat}
    \item \textbf{Construction:}
    \begin{equation}
        \text{whitecollar\_dummy} = \mathbb{1}(\texttt{subcat}=11)
    \end{equation}
    \item \textbf{Purpose:} Identifies white-collar jobs
\end{itemize}

\paragraph{bluecollar\_dummy}
\begin{itemize}
    \item \textbf{Source:} Employment subcategory \texttt{subcat}
    \item \textbf{Construction:}
    \begin{equation}
        \text{bluecollar\_dummy} = \mathbb{1}(\texttt{subcat}=12)
    \end{equation}
    \item \textbf{Purpose:} Identifies blue-collar jobs
\end{itemize}

\paragraph{othercollar\_dummy}
\begin{itemize}
    \item \textbf{Source:} Employment subcategory \texttt{subcat}
    \item \textbf{Construction:}
    \begin{equation}
        \text{othercollar\_dummy} = \mathbb{1}(\texttt{subcat} \notin \{11,12\} \text{ and not missing})
    \end{equation}
    \item \textbf{Purpose:} Captures other employment types outside the main white-/blue-collar distinction
\end{itemize}

\paragraph{subcat}
\begin{itemize}
    \item \textbf{Source:} Employment subcategory
    \item \textbf{Construction:} Collapsed job-type indicator:
    \begin{equation}
        \text{subcat} =
        \begin{cases}
            0 & \text{if original subcat} = 11\\
            1 & \text{if original subcat} \in \{12,14\}
        \end{cases}
    \end{equation}
    \item \textbf{Purpose:} Provides a compact collar/job-type distinction used in earlier model specifications and retained in V4
\end{itemize}

\paragraph{live\_birth\_event}
\begin{itemize}
    \item \textbf{Source:} Raw qualification/event data (\texttt{qualifikationen\_data})
    \item \textbf{Construction:} Indicator for a live birth event in a given person-year, based on \texttt{qu\_qual = "49"}:
    \begin{equation}
        \text{live\_birth\_event}_{i,t} = \mathbb{1}(\text{worker } i \text{ has a live birth event in year } t)
    \end{equation}
    \item \textbf{Purpose:} Captures fertility-related labor market effects
\end{itemize}

\paragraph{n\_children}
\begin{itemize}
    \item \textbf{Source:} Derived from \texttt{live\_birth\_event}
    \item \textbf{Construction:} Cumulative sum of live birth events within person:
    \begin{equation}
        \text{n\_children}_{i,t} = \sum_{\tau \leq t} \text{live\_birth\_event}_{i,\tau}
    \end{equation}
    \item \textbf{Implementation detail:} Set to zero for men in the code:
    \begin{equation}
        \text{n\_children}_{i,t} = 0 \quad \text{if} \quad \texttt{frau}=0
    \end{equation}
    \item \textbf{Purpose:} Captures family-related wage dynamics, especially around motherhood
\end{itemize}

\paragraph{contryear}
\begin{itemize}
    \item \textbf{Source:} Administrative yearly panel structure
    \item \textbf{Construction:} Calendar year of the observation
    \item \textbf{Purpose:} Captures aggregate time effects, institutional changes, and macroeconomic conditions
\end{itemize}

\paragraph{bundesland}
\begin{itemize}
    \item \textbf{Source:} Firm-level administrative data
    \item \textbf{Construction:} Region of the employer, merged by $(\texttt{benr}, \texttt{contryear})$
    \item \textbf{Purpose:} Captures regional wage differentials
\end{itemize}

\paragraph{nace08\_2}
\begin{itemize}
    \item \textbf{Source:} Firm-level administrative data
    \item \textbf{Construction:} Industry classification at the 2-digit NACE 2008 level, merged by $(\texttt{benr}, \texttt{contryear})$
    \item \textbf{Purpose:} Captures industry-specific wage structures
\end{itemize}









\section{Model Architecture and Configuration}

\subsection{Model Type: XGBoost Gradient Boosting}

\begin{itemize}
    \item \textbf{Algorithm:} Extreme Gradient Boosting (XGBoost)
    \item \textbf{Package:} \texttt{xgboost} in R
    \item \textbf{Model class:} Regression (\texttt{reg:squarederror})
    \item \textbf{Boosting type:} Tree-based gradient boosting
\end{itemize}



\begin{comment}
\subsection{Model Objective}

\textbf{Prediction task:} Predict true uncensored log daily wage (LZ) from censored BG features

\textbf{Mathematical formulation:}
\begin{equation}
    \hat{y} = f(X_{BG})
\end{equation}
where:
\begin{itemize}
    \item $\hat{y}$ = predicted log daily wage (approximates LZ true wage)
    \item $X_{BG}$ = feature vector from BG data
    \item $f$ = XGBoost model
\end{itemize}

\textbf{Loss function:} Root Mean Squared Error (RMSE)
\begin{equation}
    \text{RMSE} = \sqrt{\frac{1}{n}\sum_{i=1}^{n}(\hat{y}_i - y_i^{\text{true}})^2}
\end{equation}
where $y^{\text{true}}$ = \texttt{lz\_dwage\_log}

\end{comment}



\subsection{Hyperparameters}

The model is configured with the following hyperparameters:

\begin{table}[H]
\centering
\begin{tabular}{lll}
\toprule
\textbf{Parameter} & \textbf{Value} & \textbf{Purpose} \\
\midrule
objective & reg:squarederror & Regression with MSE loss \\
eval\_metric & rmse & Evaluation metric \\
eta & 0.1 & Learning rate (shrinkage) \\
max\_depth & 6 & Maximum tree depth \\
min\_child\_weight & 1 & Min sum of instance weight in child \\
subsample & 0.8 & Fraction of samples for each tree \\
colsample\_bytree & 0.8 & Fraction of features for each tree \\
reg\_alpha & 0.1 & L1 regularization \\
reg\_lambda & 1 & L2 regularization \\
seed & 123 & Reproducibility \\
\bottomrule
\end{tabular}
\caption{XGBoost hyperparameters and their purposes}
\label{tab:hyperparameters}
\end{table}

\textbf{Hyperparameter explanations:}
\begin{itemize}
    \item \textbf{eta (0.1):} Conservative learning rate reduces overfitting
    \item \textbf{max\_depth (6):} Moderately deep trees capture interactions without overfitting
    \item \textbf{subsample (0.8):} Prevents overfitting through random sampling
    \item \textbf{colsample\_bytree (0.8):} Feature bagging for robustness
    \item \textbf{reg\_alpha \& reg\_lambda:} Regularization to prevent overfitting
\end{itemize}

\subsection{Training Process}

\subsubsection{Cross-validation for optimal rounds}

The optimal number of boosting rounds is determined using 5-fold cross-validation with early stopping:

\begin{lstlisting}[language=R]
xgb_cv <- xgb.cv(
  params = xgb_params,
  data = dtrain,
  nrounds = 1000,              # Maximum iterations
  nfold = 5,                   # 5-fold cross-validation
  early_stopping_rounds = 50,  # Stop if no improvement
  verbose = FALSE
)

best_nrounds <- xgb_cv$best_iteration
\end{lstlisting}

\subsubsection{Final model training}

\begin{lstlisting}[language=R]
xgb_model <- xgb.train(
  params = xgb_params,
  data = dtrain,
  nrounds = best_nrounds,
  watchlist = list(train = dtrain, test = dtest),
  verbose = FALSE
)
\end{lstlisting}

\subsection{Train/Test Split}

\begin{itemize}
    \item \textbf{Split ratio:} 80\% training, 20\% testing
    \item \textbf{Random seed:} 123 (for reproducibility)
\end{itemize}

\section{Estimation Strategy}

\subsection{Single Pooled Model Approach}

\textbf{Unlike Tobit model} (which estimates separate models by gender $\times$ age group $\times$ year):
\begin{itemize}
    \item \textbf{XGBoost uses one pooled model} across all observations
    \item Gender and age are included as features, not stratification variables
    \item Model should learn interactions automatically through the tree structure
\end{itemize}

\begin{comment}
\textbf{Advantages:}
\begin{itemize}
    \item Borrows strength across groups
    \item Learns complex interactions % (e.g., gender $\times$ age $\times$ firm size)
    \item Captures non-linear patterns
\end{itemize}
\end{comment}


\subsection{Training vs Production Use}

\textbf{Phase 1: Training phase}

For observations with both BG and LZ data:
\begin{itemize}
    \item \textbf{Input:} BG features (censored wages, firm info, individual history)
    \item \textbf{Target:} LZ true wages (uncensored)
    \item \textbf{Model learns:} How to recover true wages from censored BG patterns
\end{itemize}

\noindent\textbf{Phase 2: Production use} 

For observations with only BG data (no LZ):
\begin{itemize}
    \item \textbf{Input:} Same features and data from BG 
    \item \textbf{Output:} Predicted true wages (de-censored estimates)
    \item The model returns an estimate of what the real wage would be based on the LZ wage
\end{itemize}

\subsection{Key Insight: Censoring Pattern Learning}



\begin{comment}
The model learns that:
\begin{enumerate}
    \item \textbf{Censored observations} (\texttt{interval\_upper} $>$ \texttt{interval\_lower}) have true wages above threshold
    \item \textbf{Individual history} (\texttt{mean\_censored}, \texttt{mean\_wage}) predicts how far above
    \item \textbf{Coworker patterns} (\texttt{mean\_wage\_coworkers}, \texttt{mean\_censored\_coworkers}) indicate firm quality
    \item \textbf{Demographics} (age, gender) interact with censoring patterns
\end{enumerate} 
\end{comment}




\begin{comment}
    
\section{Model Performance Evaluation}

\subsection{Comparison Approaches}

\textbf{Three wage measures compared:}
\begin{enumerate}
    \item \textbf{BG Original:} Censored wages from Beitragsgrundlagen (current practice)
    \item \textbf{XGBoost Predicted:} Model predictions (proposed method)
    \item \textbf{LZ Uncensored:} True wages from Lohnzettel (gold standard, test set only)
\end{enumerate}

\subsection{Performance Metrics}

\textbf{Primary metrics:}
\begin{itemize}
    \item \textbf{RMSE} (Root Mean Squared Error): $\sqrt{\frac{1}{n}\sum(\text{predicted} - \text{true})^2}$
    \item \textbf{R²} (R-squared): Proportion of variance explained
    \item \textbf{MAE} (Mean Absolute Error): $\frac{1}{n}\sum|\text{predicted} - \text{true}|$
\end{itemize}

\textbf{Computed for:}
\begin{itemize}
    \item Overall (all observations)
    \item By censoring status (censored vs uncensored)
    \item Train vs test sets
\end{itemize}

\subsection{Baseline Comparison}

\textbf{Current method} (BG wages):
\begin{equation}
    \text{RMSE}_{BG} = \sqrt{\frac{1}{n}\sum(\text{BG\_wage} - \text{LZ\_true})^2}
\end{equation}

\textbf{XGBoost improvement:}
\begin{align}
    \text{Improvement} &= \text{RMSE}_{BG} - \text{RMSE}_{XGBoost}\\
    \text{Relative improvement} &= \frac{\text{RMSE}_{BG} - \text{RMSE}_{XGBoost}}{\text{RMSE}_{BG}}
\end{align}
\end{comment}


\section{Visualization and Diagnostic Plots}

\subsection{Distribution Comparisons}

\textbf{Plot 1 (P1):} Overall wage distribution
\begin{itemize}
    \item Histograms of BG Original vs XGBoost Predicted
    \item Shows if model shifts distribution appropriately
\end{itemize}

\begin{figure}[H]
\centering
\includegraphics[width=0.8\textwidth]{xgb_plot_01}
\caption{XGBoost: Distribution Comparison (BG Original vs XGBoost Predicted)}
\label{fig:xgb_plot_01}
\end{figure}

\textbf{Plot 2 (P2):} Distribution by censoring status
\begin{itemize}
    \item Left panel: Uncensored observations
    \item Right panel: Censored observations
    \item Reveals model behavior for each group
\end{itemize}

\begin{figure}[H]
\centering
\includegraphics[width=0.9\textwidth]{xgb_plot_02}
\caption{XGBoost:  Wage Distribution by Censoring Status}
\label{fig:xgb_plot_02}
\end{figure}

\subsection{Three-Way Comparisons with Gold Standard}

\textbf{Plot 11 (P11):} All three wage sources (overall)
\begin{itemize}
    \item BG Original (blue): Censored wages
    \item XGBoost Predicted (red): Model estimates
    \item LZ Uncensored (green): True wages
    \item Shows model's ability to approximate true distribution
\end{itemize}

\begin{figure}[H]
\centering
\includegraphics[width=0.9\textwidth]{xgb_plot_11}
\caption{XGBoost: BG Original vs XGBoost Predicted vs LZ Uncensored}
\label{fig:xgb_plot_11}
\end{figure}

\textbf{Plot 12 (P12):} Three-way by censoring status
\begin{itemize}
    \item Reveals if model successfully de-censors top-coded wages
    \item Shows if uncensored wages remain accurate
\end{itemize}

\begin{figure}[H]
\centering
\includegraphics[width=0.9\textwidth]{xgb_plot_12}
\caption{XGBoost: Daily Wage Distribution by Censoring Status}
\label{fig:xgb_plot_12}
\end{figure}

\textbf{Plot 13 (P13):} Density overlay
\begin{itemize}
    \item Smooth density estimates for visual comparison
    \item Easier to see distributional differences than histograms
\end{itemize}


\begin{figure}[H]
\centering
\includegraphics[width=0.9\textwidth]{xgb_plot_13}
\caption{XGBoost: Density Overlay (BG vs XGBoost vs LZ)}
\label{fig:xgb_plot_13}
\end{figure}



\newpage
\begin{table}[H]
\centering
\begin{tabular}{lrrrrrr}
\toprule
Method & n\_obs & Mean Error & Median Error & SD Error & MAE & RMSE \\
\midrule
Tobit      & 34972  & -0.0500 & -0.0442 & 0.4650 & 0.2810 & 0.4670 \\
XGBoost V1 & 6941   & -0.0031 &  0.0010 & 0.1450 & 0.0910 & 0.1450 \\
XGBoost V2 & 6941   &  0.0021 &  0.0040 & 0.1090 & 0.0666 & 0.1090 \\
XGBoost V3 & 6941   &  0.0016 &  0.0034 & 0.0995 & 0.0636 & 0.0995 \\
\bottomrule
\end{tabular}
\caption{Error summary with V3 (censored observations only). Error = Predicted - LZ true wage.}
\label{tab:error_summary_v3_censored}
\end{table}

\begin{table}[H]
\centering
\begin{tabular}{lrrrrrr}
\toprule
Method & n\_obs & Mean Error & Median Error & SD Error & MAE & RMSE \\
\midrule
Tobit      & 509143 &  0.1350 &  0.1020  & 0.2990 & 0.1770 & 0.3280 \\
XGBoost V1 & 101833 & -0.0001 & -0.0113  & 0.1580 & 0.0739 & 0.1580 \\
XGBoost V2 & 101833 &  0.0003 & -0.0095  & 0.1200 & 0.0486 & 0.1200 \\
XGBoost V3 & 101833 &  0.0004 & -0.0091  & 0.1170 & 0.0481 & 0.1170 \\
\bottomrule
\end{tabular}
\caption{Error summary with V3 (all observations). Error = Predicted - LZ true wage.}
\label{tab:error_summary_v3_all}
\end{table}

\begin{table}[H]
\centering
\begin{tabular}{lrrrrrrrr}
\toprule
Method & n\_obs & Mean Error & Median Error & SD Error & MAE & RMSE & Rank & vs Worst (\%) \\
\midrule
XGBoost V3 & 6941  &  0.0016 &  0.0034 & 0.0995 & 0.0636 & 0.0995 & 1 & 78.7 \\
XGBoost V2 & 6941  &  0.0021 &  0.0040 & 0.1090 & 0.0666 & 0.1090 & 2 & 76.7 \\
XGBoost V1 & 6941  & -0.0031 &  0.0010 & 0.1450 & 0.0910 & 0.1450 & 3 & 69.0 \\
Tobit      & 34972 & -0.0500 & -0.0442 & 0.4650 & 0.2810 & 0.4670 & 4 & 0.0 \\
\bottomrule
\end{tabular}
\caption{Final performance ranking (censored observations).}
\label{tab:ranking_v3_censored}
\end{table}

\begin{table}[H]
\centering
\begin{tabular}{lrrrrrrrr}
\toprule
Method & n\_obs & Mean Error & Median Error & SD Error & MAE & RMSE & Rank & vs Worst (\%) \\
\midrule
XGBoost V3 & 101833 &  0.0004 & -0.0091 & 0.1170 & 0.0481 & 0.1170 & 1 & 64.3 \\
XGBoost V2 & 101833 &  0.0003 & -0.0095 & 0.1200 & 0.0486 & 0.1200 & 2 & 63.3 \\
XGBoost V1 & 101833 & -0.0001 & -0.0113 & 0.1580 & 0.0739 & 0.1580 & 3 & 51.9 \\
Tobit      & 509143 &  0.1350 &  0.1020 & 0.2990 & 0.1770 & 0.3280 & 4 & 0.0 \\
\bottomrule
\end{tabular}
\caption{Final performance ranking (all observations).}
\label{tab:ranking_v3_all}
\end{table}




\newpage
\section{V4 and V2 Results from Current Sample Outputs}

This section summarizes the estimation results from the two current XGBoost output files:
\begin{itemize}
    \item \texttt{model\_v4\_v2\_sample\_output.R}
    \item \texttt{model\_v4\_v2\_sample\_output\_2.R}
\end{itemize}

The two outputs differ substantially in sample size. The first output corresponds to the larger run based on the full pre-model dataset content in the analysis window, while the second output corresponds to the much smaller run. This makes them useful for comparing how model fit changes with sample scale.

\subsection{Core Model Performance: Larger Output Sample}

For the larger run, the training sample contains 3,218,020 observations with both BG and LZ information. The corresponding test sample size is 643,604 observations.

\begin{table}[H]
\centering
\begin{tabular}{lrrrr}
\toprule
Model & Test RMSE vs LZ & Test $R^2$ vs LZ & BG Baseline Test RMSE & RMSE Reduction vs BG \\
\midrule
XGBoost V1 & 0.2561 & 0.8285 & 0.2854 & 0.0293 \\
XGBoost V2 & 0.2076 & 0.8874 & 0.2854 & 0.0778 \\
XGBoost V4 & 0.1972 & 0.8985 & 0.2854 & 0.0882 \\
\bottomrule
\end{tabular}
\caption{Model performance for the larger output sample.}
\label{tab:v4_large_performance}
\end{table}

\noindent
The main comparison within this run is between V2 and V4:
\begin{itemize}
    \item V2 Test RMSE: 0.2076
    \item V4 Test RMSE: 0.1972
    \item RMSE reduction of V4 over V2: 0.0104
    \item $R^2$ improvement of V4 over V2: 0.0111
\end{itemize}

\subsection{Core Model Performance: Smaller Output Sample}

For the smaller run, the training sample contains 313,981 observations with both BG and LZ information. The corresponding test sample size is 62,797 observations.

\begin{table}[H]
\centering
\begin{tabular}{lrrrr}
\toprule
Model & Test RMSE vs LZ & Test $R^2$ vs LZ & BG Baseline Test RMSE & RMSE Reduction vs BG \\
\midrule
XGBoost V1 & 0.1147 & 0.9645 & 0.2623 & 0.1475 \\
XGBoost V2 & 0.0843 & 0.9808 & 0.2623 & 0.1780 \\
XGBoost V4 & 0.0642 & 0.9888 & 0.2623 & 0.1981 \\
\bottomrule
\end{tabular}
\caption{Model performance for the smaller output sample.}
\label{tab:v4_small_performance}
\end{table}

\noindent
The main comparison within this run is again between V2 and V4:
\begin{itemize}
    \item V2 Test RMSE: 0.0843
    \item V4 Test RMSE: 0.0642
    \item RMSE reduction of V4 over V2: 0.0201
    \item $R^2$ improvement of V4 over V2: 0.0080
\end{itemize}

\subsection{Comparative Error Summary: Larger Output Sample}

\begin{table}[H]
\centering
\begin{tabular}{lrrrrrr}
\toprule
Method & n\_obs & Mean Error & Median Error & SD Error & MAE & RMSE \\
\midrule
XGBoost V4 & 643604 & 0.0004 & -0.0120 & 0.1970 & 0.0671 & 0.1970 \\
XGBoost V2 & 643604 & 0.0003 & -0.0155 & 0.2080 & 0.0706 & 0.2080 \\
XGBoost V1 & 643604 & 0.0004 & -0.0157 & 0.2560 & 0.1220 & 0.2560 \\
Tobit      & 3217965 & 0.1760 & 0.1350 & 0.3130 & 0.2140 & 0.3590 \\
\bottomrule
\end{tabular}
\caption{Error summary against LZ true wages, all observations, larger output sample.}
\label{tab:error_summary_large_all}
\end{table}

\begin{table}[H]
\centering
\begin{tabular}{lrrrrrr}
\toprule
Method & n\_obs & Mean Error & Median Error & SD Error & MAE & RMSE \\
\midrule
XGBoost V4 & 46235 & 0.0020 & 0.0131 & 0.2190 & 0.1260 & 0.2190 \\
XGBoost V2 & 46235 & 0.0023 & 0.0165 & 0.2380 & 0.1400 & 0.2380 \\
XGBoost V1 & 46235 & -0.0143 & 0.0014 & 0.2820 & 0.1720 & 0.2820 \\
Tobit      & 231101 & -0.0218 & -0.0118 & 0.4730 & 0.2820 & 0.4730 \\
\bottomrule
\end{tabular}
\caption{Error summary against LZ true wages, censored observations only, larger output sample.}
\label{tab:error_summary_large_censored}
\end{table}

\subsection{Comparative Error Summary: Smaller Output Sample}

\begin{table}[H]
\centering
\begin{tabular}{lrrrrrr}
\toprule
Method & n\_obs & Mean Error & Median Error & SD Error & MAE & RMSE \\
\midrule
XGBoost V4 & 62797 & 0.0006 & -0.0029 & 0.0642 & 0.0315 & 0.0642 \\
XGBoost V2 & 62797 & 0.0008 & -0.0059 & 0.0843 & 0.0385 & 0.0843 \\
XGBoost V1 & 62797 & 0.0004 & -0.0082 & 0.1150 & 0.0569 & 0.1150 \\
Tobit      & 313970 & 0.1270 & 0.0939 & 0.2760 & 0.1650 & 0.3040 \\
\bottomrule
\end{tabular}
\caption{Error summary against LZ true wages, all observations, smaller output sample.}
\label{tab:error_summary_small_all}
\end{table}

\begin{table}[H]
\centering
\begin{tabular}{lrrrrrr}
\toprule
Method & n\_obs & Mean Error & Median Error & SD Error & MAE & RMSE \\
\midrule
XGBoost V4 & 4319 & -0.0008 & 0.0015 & 0.0512 & 0.0373 & 0.0512 \\
XGBoost V2 & 4319 & -0.0017 & 0.0009 & 0.0609 & 0.0439 & 0.0609 \\
XGBoost V1 & 4319 & -0.0047 & -0.0003 & 0.0872 & 0.0615 & 0.0873 \\
Tobit      & 21486 & -0.0002 & 0.0002 & 0.3910 & 0.2640 & 0.3910 \\
\bottomrule
\end{tabular}
\caption{Error summary against LZ true wages, censored observations only, smaller output sample.}
\label{tab:error_summary_small_censored}
\end{table}

\subsection{Comprehensive RMSE Tables from the Two Outputs}

\subsubsection{Larger Output Sample}

\begin{table}[H]
\centering
\small
\begin{tabular}{llrrrrr}
\toprule
Method & Observation Type & n\_obs & RMSE & MAE & Mean Error & SD Error \\
\midrule
Tobit & All (BG target) & 3217965 & 0.2520 & 0.1830 & 0.1830 & 0.1740 \\
Tobit & All (LZ target) & 3217965 & 0.3590 & 0.2140 & 0.1760 & 0.3130 \\
XGBoost V1 & All & 643604 & 0.2560 & 0.1220 & 0.0004 & 0.2560 \\
XGBoost V2 & All & 643604 & 0.2080 & 0.0706 & 0.0003 & 0.2080 \\
XGBoost V4 & All & 643604 & 0.1970 & 0.0671 & 0.0004 & 0.1970 \\
Tobit & Censored (BG target) & 231101 & 0.3720 & 0.3050 & 0.3050 & 0.2130 \\
Tobit & Censored (LZ target) & 231101 & 0.4730 & 0.2820 & -0.0218 & 0.4730 \\
XGBoost V1 & Censored & 46235 & 0.2820 & 0.1720 & -0.0143 & 0.2820 \\
XGBoost V2 & Censored & 46235 & 0.2380 & 0.1400 & 0.0023 & 0.2380 \\
XGBoost V4 & Censored & 46235 & 0.2190 & 0.1260 & 0.0020 & 0.2190 \\
Tobit & Uncensored (BG target) & 2986864 & 0.2400 & 0.1730 & 0.1730 & 0.1660 \\
Tobit & Uncensored (LZ target) & 2986864 & 0.3490 & 0.2090 & 0.1910 & 0.2920 \\
XGBoost V1 & Uncensored & 597369 & 0.2540 & 0.1180 & 0.0015 & 0.2540 \\
XGBoost V2 & Uncensored & 597369 & 0.2050 & 0.0653 & 0.0001 & 0.2050 \\
XGBoost V4 & Uncensored & 597369 & 0.1950 & 0.0625 & 0.0003 & 0.1950 \\
\bottomrule
\end{tabular}
\caption{Comprehensive RMSE comparison, larger output sample.}
\label{tab:comprehensive_rmse_large}
\end{table}

\subsubsection{Smaller Output Sample}

\begin{table}[H]
\centering
\small
\begin{tabular}{llrrrrr}
\toprule
Method & Observation Type & n\_obs & RMSE & MAE & Mean Error & SD Error \\
\midrule
Tobit & All (BG target) & 313970 & 0.1900 & 0.1320 & 0.1320 & 0.1370 \\
Tobit & All (LZ target) & 313970 & 0.3040 & 0.1650 & 0.1270 & 0.2760 \\
XGBoost V1 & All & 62797 & 0.1150 & 0.0569 & 0.0004 & 0.1150 \\
XGBoost V2 & All & 62797 & 0.0843 & 0.0385 & 0.0008 & 0.0843 \\
XGBoost V4 & All & 62797 & 0.0642 & 0.0315 & 0.0006 & 0.0642 \\
Tobit & Censored (BG target) & 21486 & 0.3930 & 0.3200 & 0.3200 & 0.2300 \\
Tobit & Censored (LZ target) & 21486 & 0.3910 & 0.2640 & -0.0002 & 0.3910 \\
XGBoost V1 & Censored & 4319 & 0.0873 & 0.0615 & -0.0047 & 0.0872 \\
XGBoost V2 & Censored & 4319 & 0.0609 & 0.0439 & -0.0017 & 0.0609 \\
XGBoost V4 & Censored & 4319 & 0.0512 & 0.0373 & -0.0008 & 0.0512 \\
Tobit & Uncensored (BG target) & 292484 & 0.1660 & 0.1180 & 0.1180 & 0.1160 \\
Tobit & Uncensored (LZ target) & 292484 & 0.2960 & 0.1570 & 0.1360 & 0.2630 \\
XGBoost V1 & Uncensored & 58478 & 0.1160 & 0.0566 & 0.0008 & 0.1160 \\
XGBoost V2 & Uncensored & 58478 & 0.0858 & 0.0381 & 0.0010 & 0.0858 \\
XGBoost V4 & Uncensored & 58478 & 0.0651 & 0.0311 & 0.0007 & 0.0651 \\
\bottomrule
\end{tabular}
\caption{Comprehensive RMSE comparison, smaller output sample.}
\label{tab:comprehensive_rmse_small}
\end{table}

\subsection{Variables Included in Each XGBoost Model}

\subsubsection{XGBoost V1 Variables}

\begin{itemize}
    \item \texttt{age}
    \item \texttt{mean\_wage}
    \item \texttt{mean\_censored}
    \item \texttt{firm\_size}
    \item \texttt{firm\_size\_square}
    \item \texttt{oneyear\_dummy}
    \item \texttt{oneperson\_dummy}
    \item \texttt{smallfirm\_dummy}
    \item \texttt{mean\_wage\_coworkers}
    \item \texttt{mean\_censored\_coworkers}
\end{itemize}

\subsubsection{XGBoost V2 Variables}

\begin{itemize}
    \item \texttt{age}
    \item \texttt{mean\_wage}
    \item \texttt{mean\_censored}
    \item \texttt{firm\_size}
    \item \texttt{firm\_size\_square}
    \item \texttt{oneyear\_dummy}
    \item \texttt{oneperson\_dummy}
    \item \texttt{smallfirm\_dummy}
    \item \texttt{mean\_wage\_coworkers}
    \item \texttt{mean\_censored\_coworkers}
    \item \texttt{frau}
    \item \texttt{birthyear}
    \item \texttt{num\_years}
    \item \texttt{interval\_lower}
    \item \texttt{interval\_upper}
    \item \texttt{bg\_censored\_flag}
\end{itemize}

\subsubsection{XGBoost V4 Variables}

\begin{itemize}
    \item \texttt{age}
    \item \texttt{exper}
    \item \texttt{tenure}
    \item \texttt{mean\_wage}
    \item \texttt{mean\_censored}
    \item \texttt{firm\_size}
    \item \texttt{firm\_size\_square}
    \item \texttt{oneyear\_dummy}
    \item \texttt{oneperson\_dummy}
    \item \texttt{smallfirm\_dummy}
    \item \texttt{mean\_wage\_coworkers}
    \item \texttt{mean\_censored\_coworkers}
    \item \texttt{interval\_lower}
    \item \texttt{interval\_upper}
    \item \texttt{bg\_censored\_flag}
    \item \texttt{frau}
    \item \texttt{stsl\_missing}
    \item \texttt{austrian\_dummy}
    \item \texttt{germany\_dummy}
    \item \texttt{eu\_dummy}
    \item \texttt{educ\_cat}
    \item \texttt{educ\_missing}
    \item \texttt{educ\_imputed}
    \item \texttt{whitecollar\_dummy}
    \item \texttt{bluecollar\_dummy}
    \item \texttt{othercollar\_dummy}
    \item \texttt{subcat}
    \item \texttt{juj}
    \item \texttt{return}
    \item \texttt{year\_entry}
    \item \texttt{age\_entry}
    \item \texttt{num\_years}
    \item \texttt{live\_birth\_event}
    \item \texttt{n\_children}
    \item \texttt{contryear}
    \item \texttt{bundesland}
    \item \texttt{nace08\_2}
\end{itemize}


\begin{comment}
   \subsection{Interpretation of the New V4 Results}

Across both output files, the ranking of models is stable:
\begin{enumerate}
    \item XGBoost V4 performs best
    \item XGBoost V2 performs second best
    \item XGBoost V1 performs third
    \item Tobit performs worst in prediction against LZ true wages
\end{enumerate}

The improvement from V2 to V4 is present in both runs, but the gain is noticeably larger in the smaller output sample than in the larger output sample. This is especially visible in the test RMSE reduction:
\begin{itemize}
    \item Smaller output sample: V4 improves on V2 by 0.0201 RMSE
    \item Larger output sample: V4 improves on V2 by 0.0104 RMSE
\end{itemize}

This pattern is consistent with the concern that the model's apparent fit deteriorates as sample size increases. Even so, V4 remains the best-performing specification in both samples, and it substantially outperforms Tobit on both all observations and censored observations only.

 
\end{comment}


\begin{comment}

\subsection{Interpretation Guidelines}

\textbf{Success indicators:}
\begin{itemize}
    \item XGBoost distribution closer to LZ than BG (especially in censored group)
    \item Right tail of XGBoost extends beyond BG censoring threshold
    \item Uncensored XGBoost predictions match BG and LZ closely
\end{itemize}

\textbf{Potential issues:}
\begin{itemize}
    \item Over-prediction: XGBoost consistently above LZ
    \item Under-prediction: XGBoost doesn't extend far enough into right tail
    \item Variance mismatch: Too narrow or too wide distribution
\end{itemize}

\section{Feature Importance Analysis}

\subsection{XGBoost Feature Importance Metrics}

\textbf{Gain:} Total improvement in loss function from splits on feature
\begin{itemize}
    \item Higher gain = more important for prediction accuracy
    \item Reports: \texttt{importance\_matrix\_v2}
\end{itemize}

\textbf{Interpretation:}
\begin{itemize}
    \item Top features drive most of model's predictive power
    \item Can identify which BG characteristics best predict true wages
    \item Compare to economic intuition (e.g., should coworker wages matter?)
\end{itemize}

\subsection{Expected Important Features}

\textbf{Hypothesized ranking:}
\begin{enumerate}
    \item \texttt{interval\_lower} / \texttt{interval\_upper}: Direct wage information
    \item \texttt{mean\_wage}: Individual fixed effect
    \item \texttt{mean\_wage\_coworkers}: Firm wage premium
    \item \texttt{bg\_censored\_flag}: Indicates top-coding
    \item \texttt{age}: Life-cycle profile
    \item \texttt{firm\_size}: Firm size premium
\end{enumerate}
\end{comment}


\begin{comment}
\section{Model Versions and Extensions}

\subsection{Version Comparison}

\textbf{V1} (Original feature set):
\begin{itemize}
    \item Fewer features, faster training
    \item Baseline performance
\end{itemize}

\textbf{V2} (Expanded feature set):
\begin{itemize}
    \item Includes all Tobit features + additional BG features
    \item Tests if more information improves predictions
    \item Reports relative improvement: V2\_RMSE - V1\_RMSE
\end{itemize}

\subsection{Future Extensions}

\textbf{Potential improvements:}
\begin{enumerate}
    \item \textbf{Industry controls:} Add NACE codes
    \item \textbf{Regional controls:} Add Bundesland
    \item \textbf{Temporal features:} Year trends, economic conditions
    \item \textbf{Interaction terms:} Manual feature engineering
    \item \textbf{Ensemble methods:} Combine XGBoost with Tobit predictions
\end{enumerate}

\section{Production Deployment}

\subsection{Prediction Function}

\begin{lstlisting}[language=R]
predict_true_wage <- function(bg_data, 
                              model = xgb_model_v2, 
                              feature_columns = feature_cols_v2) {
  
  # Extract and prepare features
  prediction_features <- bg_data %>%
    select(all_of(feature_columns)) %>%
    mutate(across(everything(), 
           ~ifelse(is.na(.) | is.infinite(.), 0, .))) %>%
    as.matrix()
  
  # Predict
  predicted_wages <- predict(model, 
                            xgb.DMatrix(data = prediction_features))
  
  return(predicted_wages)
}
\end{lstlisting}

\subsection{Application to Future Data}

\textbf{Requirements:}
\begin{enumerate}
    \item New BG data with same variable construction
    \item All feature variables must be computed identically
    \item CPI adjustment to 2000 Euros
    \item Log transformation
\end{enumerate}

\textbf{Output:} De-censored log daily wage estimates

\subsection{Limitations and Assumptions}

\textbf{Key assumptions:}
\begin{enumerate}
    \item Relationship between BG features and LZ wages remains stable over time
    \item Censoring rules haven't changed substantially
    \item Sample representativeness: Training sample generalizes to full population
\end{enumerate}

\textbf{Limitations:}
\begin{enumerate}
    \item Predictions only as good as training data overlap (BG-LZ match)
    \item Cannot predict beyond highest observed LZ wages
    \item Model uncertainty not quantified (point predictions only)
    \item Requires re-training if censoring thresholds change significantly
\end{enumerate}

\section{Comparison with Tobit Approach}

\subsection{Methodological Differences}

\begin{table}[H]
\centering
\small
\begin{tabular}{p{3cm}p{5cm}p{5cm}}
\toprule
\textbf{Aspect} & \textbf{Tobit Model} & \textbf{XGBoost Model} \\
\midrule
Estimation & Separate models per gender$\times$age$\times$year & Single pooled model \\
Functional form & Linear (log-linear) & Non-linear (tree-based) \\
Interactions & Must specify manually & Learned automatically \\
Censoring & Explicitly modeled intervals & Learned from patterns \\
Interpretability & High (coefficients) & Lower (black box) \\
Flexibility & Limited to specified form & Highly flexible \\
Sample size needs & Requires sufficient obs per group & Leverages full sample \\
Extrapolation & Based on linear trends & Based on tree splits \\
\bottomrule
\end{tabular}
\caption{Comparison of Tobit and XGBoost approaches}
\label{tab:comparison}
\end{table}

\subsection{Expected Performance Differences}

\textbf{XGBoost advantages:}
\begin{itemize}
    \item Captures non-linearities in age-wage profiles
    \item Learns complex interactions (e.g., gender $\times$ firm size $\times$ censoring)
    \item No need to specify functional form
    \item Robust to outliers
\end{itemize}

\textbf{Tobit advantages:}
\begin{itemize}
    \item Economic interpretability
    \item Explicit uncertainty quantification
    \item Well-understood statistical properties
    \item Handles censoring mechanism directly
\end{itemize}

\subsection{Complementary Use}

\textbf{Recommendation:} Use both approaches
\begin{itemize}
    \item Tobit for interpretation and statistical inference
    \item XGBoost for prediction accuracy
    \item Compare predictions as robustness check
    \item Ensemble: Average predictions for final estimates
\end{itemize}

\section{Summary Statistics and Results}

\subsection{Sample Composition}

\textbf{To be filled with actual results:}
\begin{itemize}
    \item Total observations in training data: [N]
    \item Training set: [N\_train] (80\%)
    \item Test set: [N\_test] (20\%)
    \item Censored observations: [N\_censored] ([pct]\%)
    \item Uncensored observations: [N\_uncensored] ([pct]\%)
\end{itemize}

\subsection{Model Performance Summary}

\textbf{Test Set Performance:}

\begin{table}[H]
\centering
\begin{tabular}{lcccc}
\toprule
\textbf{Method} & \textbf{RMSE} & \textbf{R²} & \textbf{MAE} & \textbf{\% Better than BG} \\
\midrule
BG Original & [X] & [X] & [X] & --- \\
XGBoost V2 & [X] & [X] & [X] & [X]\% \\
Improvement & [X] & [X] & [X] & --- \\
\bottomrule
\end{tabular}
\caption{Model performance on test set}
\label{tab:performance}
\end{table}

\textbf{By Censoring Status:}

\begin{table}[H]
\centering
\begin{tabular}{lccc}
\toprule
\textbf{Group} & \textbf{BG RMSE} & \textbf{XGB RMSE} & \textbf{Improvement} \\
\midrule
Uncensored & [X] & [X] & [X] \\
Censored & [X] & [X] & [X] \\
\bottomrule
\end{tabular}
\caption{Performance by censoring status}
\label{tab:censoring_performance}
\end{table}

\subsection{Distribution Statistics}

\textbf{Wage Distributions (Log Daily Wage, 2000 Euros):}

\begin{table}[H]
\centering
\begin{tabular}{lcccc}
\toprule
\textbf{Measure} & \textbf{BG Original} & \textbf{XGB Predicted} & \textbf{LZ True} & \textbf{Diff (XGB - LZ)} \\
\midrule
Mean & [X] & [X] & [X] & [X] \\
Median & [X] & [X] & [X] & [X] \\
SD & [X] & [X] & [X] & [X] \\
P10 & [X] & [X] & [X] & [X] \\
P90 & [X] & [X] & [X] & [X] \\
P95 & [X] & [X] & [X] & [X] \\
P99 & [X] & [X] & [X] & [X] \\
\bottomrule
\end{tabular}
\caption{Wage distribution statistics}
\label{tab:wage_distribution}
\end{table}

\section{Conclusions and Recommendations}

\subsection{Key Findings}

\begin{enumerate}
    \item \textbf{Model successfully de-censors wages:} [Describe if right tail extends appropriately]
    \item \textbf{Performance gain over BG:} [X]\% RMSE reduction
    \item \textbf{Most important features:} [List top 5 from feature importance]
    \item \textbf{Censored vs uncensored performance:} [Compare accuracy across groups]
\end{enumerate}

\subsection{Practical Implications}

\textbf{For wage research:}
\begin{itemize}
    \item XGBoost predictions can substitute for unavailable LZ data
    \item Reduces bias from censoring in wage inequality studies
    \item Enables analysis of high-wage workers
\end{itemize}

\textbf{For policy analysis:}
\begin{itemize}
    \item More accurate wage estimates for top earners
    \item Better understanding of high-wage employment
    \item Improved estimation of wage distributions
\end{itemize}

\subsection{Next Steps}

\begin{enumerate}
    \item \textbf{Validate on additional years} (2013+)
    \item \textbf{Incorporate industry and regional controls}
    \item \textbf{Develop uncertainty quantification} (prediction intervals)
    \item \textbf{Compare with other ML methods} (Random Forest, Neural Networks)
    \item \textbf{Create ensemble with Tobit predictions}
\end{enumerate}

\appendix

\section{Code Structure}

\subsection{Main Workflow}

\begin{lstlisting}[language=R]
# 1. Data preparation
training_data <- tobit_data %>%
  left_join(lohnzettel) %>%
  filter(!is.na(lz_dwage_log))

# 2. Feature matrix
X_features <- training_data %>%
  select(all_of(feature_cols)) %>%
  as.matrix()

y_target <- training_data$lz_dwage_log

# 3. Train/test split
train_indices <- sample(...)
X_train <- X_features[train_indices, ]
y_train <- y_target[train_indices]

# 4. Cross-validation
xgb_cv <- xgb.cv(...)
best_nrounds <- xgb_cv$best_iteration

# 5. Final model
xgb_model <- xgb.train(
  params = xgb_params,
  nrounds = best_nrounds
)

# 6. Predictions
test_pred <- predict(xgb_model, dtest)

# 7. Evaluation
rmse <- sqrt(mean((test_pred - y_test)^2))

# 8. Plots
# Generate P1, P2, P11, P12, P13
\end{lstlisting}

\subsection{Required Packages}

\begin{lstlisting}[language=R]
library(xgboost)    # Model training
library(dplyr)      # Data manipulation  
library(tidyr)      # Data reshaping
library(ggplot2)    # Visualization
\end{lstlisting}





\end{comment}
\end{document}